%! Equivalent Representations of Trigonometric Functions — Unit 3, Lesson 3.12 — Guided Notes (Answer Key)
\documentclass[10pt]{article}
\usepackage{precalculus-article}
\usepackage{precalculus-key}   % pulls in -boxes; do NOT also load -boxes
\newcommand{\TallMath}[1]{$\displaystyle #1\rule[-1.4em]{0pt}{3.2em}$}

\begin{document}

\pageheader{Unit 3, Lesson 3.12}{Guided Notes --- Answer Key}
\namedateperiod

\vspace{0.10in}

%----------------------------------------------------------------------
% Objectives
%----------------------------------------------------------------------
\begin{objectivebox}
\textbf{I will be able to\ldots}
\begin{enumerate}[leftmargin=1.5em, itemsep=4pt]
  \item \textbf{LO 3.12.A} --- Rewrite trig expressions using the Pythagorean identity
    $\sin^2\theta + \cos^2\theta = 1$ and its algebraic variants.
    \ansline{Use the identity to substitute one trig function for another.}
  \item \textbf{LO 3.12.B} --- Use the sum, difference, and double-angle identities for
    sine and cosine to rewrite expressions and find exact values.
    \ansline{Apply sum/difference/double-angle formulas as needed.}
  \item \textbf{LO 3.12.C} --- Solve trig equations by substituting an equivalent form to
    reduce the equation to a single trig function.
    \ansline{Convert to one function, factor or use quadratic formula, find all solutions in the given interval.}
\end{enumerate}
\end{objectivebox}

\vspace{0.08in}

%----------------------------------------------------------------------
% Vocabulary
%----------------------------------------------------------------------
\begin{vocabbox}
\begin{description}[leftmargin=0pt, itemsep=6pt]
  \item[\termblanklong{Pythagorean identity}]
    $\sin^2\theta + \cos^2\theta = $ \ans{1}.
    Derived from the unit-circle point $(\cos\theta,\sin\theta)$ and the
    \ans{Pythagorean Theorem}.

  \item[\termblanklong{sum identity (sine)}]
    $\sin(\alpha + \beta) = \sin\alpha\cdot$ \ans{$\cos\beta$}
    $+\;\cos\alpha\cdot$ \ans{$\sin\beta$}

  \item[\termblanklong{sum identity (cosine)}]
    $\cos(\alpha + \beta) = \cos\alpha\cdot$ \ans{$\cos\beta$}
    $-\;\sin\alpha\cdot$ \ans{$\sin\beta$}

  \item[\termblanklong{double-angle identity}]
    Obtained by setting $\alpha = \beta = \theta$ in the sum identity.\\
    $\sin 2\theta = $ \ans{$2\sin\theta\cos\theta$} \quad
    $\cos 2\theta = $ \ans{$\cos^2\theta - \sin^2\theta$}

  \item[\termblanklong{verify an identity}]
    Show two expressions are equal for all domain values by transforming
    \ans{one} side algebraically --- \textbf{never} add/multiply both sides.
\end{description}
\end{vocabbox}

\vspace{0.08in}

%----------------------------------------------------------------------
% Hook
%----------------------------------------------------------------------
\begin{hookbox}
\textbf{Entry Question:} Consider the equation $2\sin^2\theta + \cos\theta = 1$.

\vspace{4pt}
\textbf{Q1.} Why is it difficult to solve this equation as written?

\ansline{It has two different trig functions (sine and cosine), so we cannot factor or use the quadratic formula directly.}

\textbf{Q2.} Using $\sin^2\theta + \cos^2\theta = 1$, rewrite $\sin^2\theta$
in terms of $\cos\theta$:

\vspace{3pt}
$\sin^2\theta = $ \ans{$1 - \cos^2\theta$}

\vspace{4pt}
\textbf{Q3.} Substitute your answer into $2\sin^2\theta + \cos\theta = 1$.
What kind of equation do you get?

\ansline{$2(1-\cos^2\theta)+\cos\theta-1=0$, which is a quadratic equation in $\cos\theta$.}
\end{hookbox}

\vspace{0.08in}

%----------------------------------------------------------------------
% LO 3.12.A — Pythagorean Identity
%----------------------------------------------------------------------
\begin{notesbox}{LO 3.12.A --- The Pythagorean Identity}

\textbf{Derivation from the Unit Circle:}

A point on the unit circle has coordinates $(\cos\theta,\,\sin\theta)$.
Its distance from the origin equals the radius $= $ \ans{1}.

By the Pythagorean Theorem: $(\cos\theta)^2 + (\sin\theta)^2 = $ \ans{1}

\vspace{6pt}
\textbf{Identity 1 (standard form):}
$$\sin^2\theta + \cos^2\theta = \ans{1}$$

\vspace{4pt}
\textbf{Identity 2 (divide both sides by $\cos^2\theta$):}

\vspace{4pt}
\renewcommand{\arraystretch}{2.2}
\begin{tabular}{@{}l@{\quad}l@{}}
$\dfrac{\sin^2\theta}{\cos^2\theta} + \dfrac{\cos^2\theta}{\cos^2\theta}
= \dfrac{1}{\cos^2\theta}$
&
$\Longrightarrow\quad$ \ans{$\tan^2\theta$} $+ 1 = $ \ans{$\sec^2\theta$} \\
\end{tabular}

\vspace{4pt}
\textbf{Identity 3 (divide both sides by $\sin^2\theta$):}

\vspace{4pt}
\begin{tabular}{@{}l@{\quad}l@{}}
$\dfrac{\sin^2\theta}{\sin^2\theta} + \dfrac{\cos^2\theta}{\sin^2\theta}
= \dfrac{1}{\sin^2\theta}$
&
$\Longrightarrow\quad 1 + $ \ans{$\cot^2\theta$} $= $ \ans{$\csc^2\theta$} \\
\end{tabular}

\vspace{8pt}
\textbf{Worked Example:} If $\cos\theta = \dfrac{3}{5}$ and $\theta$ is in Quadrant IV, find $\sin\theta$.

\vspace{4pt}
$\sin^2\theta = 1 - \cos^2\theta = 1 - \left(\dfrac{3}{5}\right)^2
= 1 - $ \ans{$\dfrac{9}{25}$} $= $ \ans{$\dfrac{16}{25}$}

\vspace{4pt}
$\sin\theta = \pm$ \ans{$\dfrac{4}{5}$};\quad since $\theta$ is in Quadrant IV, $\sin\theta = $ \ans{$-\dfrac{4}{5}$}

\end{notesbox}

\vspace{0.08in}

%----------------------------------------------------------------------
% LO 3.12.B — Sum and Difference Identities
%----------------------------------------------------------------------
\begin{notesbox}{LO 3.12.B --- Sum and Difference Identities}

\textbf{Sum Identities:}
\begin{align*}
\sin(\alpha + \beta) &= \sin\alpha\cos\beta + \cos\alpha\sin\beta \\
\cos(\alpha + \beta) &= \cos\alpha\cos\beta - \sin\alpha\sin\beta
\end{align*}

\textbf{Difference Identities} (replace $\beta$ with $-\beta$):
\begin{align*}
\sin(\alpha - \beta) &= \sin\alpha\cos\beta - \cos\alpha\sin\beta \\
\cos(\alpha - \beta) &= \cos\alpha\cos\beta + \sin\alpha\sin\beta
\end{align*}

\vspace{4pt}
\textbf{Worked Example --- Find $\sin 75^\circ$ exactly.}

Write $75^\circ = 45^\circ + 30^\circ$:

\vspace{4pt}
$\sin(45^\circ + 30^\circ)
= \sin 45^\circ \cdot \cos 30^\circ + \cos 45^\circ \cdot \sin 30^\circ$

\vspace{4pt}
$= \dfrac{\sqrt{2}}{2}\cdot$ \ans{$\dfrac{\sqrt{3}}{2}$}
$+\; \dfrac{\sqrt{2}}{2}\cdot$ \ans{$\dfrac{1}{2}$}

\vspace{4pt}
$= $ \ans{$\dfrac{\sqrt{6}}{4}$} $+$ \ans{$\dfrac{\sqrt{2}}{4}$} $= $ \ans{$\dfrac{\sqrt{6}+\sqrt{2}}{4}$}

\vspace{8pt}
\textbf{Practice:} Find $\cos 15^\circ = \cos(45^\circ - 30^\circ)$ exactly.

\vspace{4pt}
$\cos(45^\circ - 30^\circ)
= \cos 45^\circ \cos 30^\circ + \sin 45^\circ \sin 30^\circ$

\vspace{4pt}
$=$ \ans{$\dfrac{\sqrt{2}}{2}$} $\cdot$ \ans{$\dfrac{\sqrt{3}}{2}$}
$+$ \ans{$\dfrac{\sqrt{2}}{2}$} $\cdot$ \ans{$\dfrac{1}{2}$}
$= $ \ans{$\dfrac{\sqrt{6}+\sqrt{2}}{4}$}

\begin{teachernote}
Note that $\cos 15^\circ = \sin 75^\circ$; this is a good observation to prompt
(cofunctions of complementary angles).
\end{teachernote}

\end{notesbox}

\vspace{0.08in}

%----------------------------------------------------------------------
% LO 3.12.B — Double-Angle Identities
%----------------------------------------------------------------------
\begin{notesbox}{LO 3.12.B --- Double-Angle Identities}

\textbf{Derivation:} Set $\alpha = \beta = \theta$ in the sum identities.

\vspace{4pt}
$\sin(2\theta) = \sin\theta\cos\theta + \cos\theta\sin\theta
= $ \ans{$2\sin\theta\cos\theta$}

\vspace{6pt}
$\cos(2\theta) = \cos\theta\cos\theta - \sin\theta\sin\theta
= $ \ans{$\cos^2\theta - \sin^2\theta$}

\vspace{8pt}
\textbf{Alternate forms of $\cos 2\theta$} (using $\sin^2\theta + \cos^2\theta = 1$):

\vspace{4pt}
$\cos 2\theta = \cos^2\theta - \sin^2\theta
= \cos^2\theta - (1 - \cos^2\theta) = $ \ans{$2\cos^2\theta - 1$}

\vspace{4pt}
$\cos 2\theta = \cos^2\theta - \sin^2\theta
= (1 - \sin^2\theta) - \sin^2\theta = $ \ans{$1 - 2\sin^2\theta$}

\vspace{8pt}
\textbf{Rearrangements for solving equations:}

\vspace{4pt}
$\cos^2\theta = \dfrac{1 + \cos 2\theta}{2}$
\hspace{1.4cm}
$\sin^2\theta = \dfrac{1 - \cos 2\theta}{2}$

\end{notesbox}

\vspace{0.08in}

%----------------------------------------------------------------------
% LO 3.12.C — Solving with Identities
%----------------------------------------------------------------------
\begin{notesbox}{LO 3.12.C --- Solving Equations Using Identities}

\textbf{Strategy:} Use an identity to rewrite everything in terms of \emph{one}
trig function, then solve the resulting algebraic equation.

\vspace{6pt}
\textbf{Worked Example:} Solve $2\sin^2\theta + \cos\theta - 1 = 0$ for $\theta \in [0, 2\pi)$.

\vspace{4pt}
\textbf{Step 1.} Substitute $\sin^2\theta = 1 - \cos^2\theta$:
$\;2(\;$ \ans{$1 - \cos^2\theta$} $\;) + \cos\theta - 1 = 0$

\vspace{4pt}
\textbf{Step 2.} Expand and collect:
\ans{$-2\cos^2\theta + \cos\theta + 1 = 0$}

\vspace{4pt}
\textbf{Step 3.} Multiply by $-1$, then factor (let $u = \cos\theta$):
$(2u + 1)(u - 1) = 0$

\vspace{4pt}
$\cos\theta = -\dfrac{1}{2}$ \quad or \quad $\cos\theta = $ \ans{1}

\vspace{4pt}
\textbf{Step 4.} Solve each case in $[0, 2\pi)$:

\vspace{2pt}
$\cos\theta = -\tfrac{1}{2}$: \quad $\theta = $ \ans{$\dfrac{2\pi}{3}$}, \ans{$\dfrac{4\pi}{3}$}

\vspace{4pt}
$\cos\theta = 1$: \quad $\theta = $ \ans{0}

\vspace{6pt}
\textbf{Solution set:} $\theta \in \big\{$ \ans{$0,\;\dfrac{2\pi}{3},\;\dfrac{4\pi}{3}$} $\big\}$

\end{notesbox}

\vspace{0.08in}

%----------------------------------------------------------------------
% Practice
%----------------------------------------------------------------------
\begin{practicebox}

\textbf{Guided Practice 1.}\quad Use a Pythagorean identity to simplify.

\vspace{4pt}
(a)\quad $\sec^2\theta - \tan^2\theta = $ \ans{1}
\hspace{1cm}
(b)\quad $1 - \sin^2\theta = $ \ans{$\cos^2\theta$}

\vspace{0.10in}
\textbf{Guided Practice 2.}\quad Evaluate exactly using an identity.

\vspace{4pt}
(a)\quad $\sin\!\left(\dfrac{\pi}{12}\right)$
\quad(\textit{Hint:} $\dfrac{\pi}{12} = \dfrac{\pi}{4} - \dfrac{\pi}{6}$)

\vspace{4pt}
$\sin\!\left(\tfrac{\pi}{4} - \tfrac{\pi}{6}\right)
= \sin\tfrac{\pi}{4}\cos\tfrac{\pi}{6} - \cos\tfrac{\pi}{4}\sin\tfrac{\pi}{6}
= \dfrac{\sqrt{2}}{2}\cdot\dfrac{\sqrt{3}}{2} - \dfrac{\sqrt{2}}{2}\cdot\dfrac{1}{2}
= $ \ans{$\dfrac{\sqrt{6}-\sqrt{2}}{4}$}

\vspace{8pt}
(b)\quad If $\cos\theta = -\dfrac{1}{3}$, find $\cos 2\theta$.

\vspace{4pt}
$\cos 2\theta = 2\cos^2\theta - 1 = 2\!\left(\dfrac{1}{9}\right) - 1 = $ \ans{$-\dfrac{7}{9}$}

\end{practicebox}

\end{document}
